Níže jsou zkrácené formulace a šablony pro sylaby a pole deklarace; upravte podle předmětu a místních pravidel.

\bigskip

\textbf{Formulace do sylabu (kopírovat/upravit):}

\begin{quote}
\begin{PromptBlock}

Použití generativní AI je v tomto předmětu povoleno pouze se souhlasem vyučujícího a za podmínky, že při odevzdání student explicitně uvede: (1) které nástroje použil, (2) konkrétní prompty a verze, (3) vlastní úpravy a reflexi, a (4) přiloží pracovní verze/drafty.
\end{PromptBlock}

\end{quote}

\begin{quote}
\begin{PromptBlock}

Použití externích generativních nástrojů k vytvoření obsahu zadání je zakázáno. Podezření na porušení integritní politiky bude řešeno podle akademických pravidel VUT.
\end{PromptBlock}

\end{quote}

\bigskip

\textbf{Pole deklarace (návrh):}
\begin{itemize}
  \item Checkbox: \textit{Použil(a) jsem externí AI nástroj při tvorbě tohoto odevzdání.}
  \item Text (150–300 znaků): nástroj, verze, použité prompty a stručně vlastní přidaná hodnota.
  \item Povinná příloha: pracovní verze / commity / screencast či jiný doklad postupu.
\end{itemize}

\bigskip

\textbf{Hodnocení (návrh):}
\begin{itemize}
  \item Transparentnost použití AI (0–5): rozsah použití a doložení promptů/výstupů.
  \item Originalita a vlastní přínos (0–10): úpravy a kritická evaluace.
  \item Doklad procesu (0–5): drafty, commity, screencasty nebo obhajoba.
\end{itemize}

\bigskip

\textbf{Krátký checklist:}
\begin{enumerate}
  \item Nástroj a verze.
  \item Přesný prompt/postup.
  \item Popis úprav a vlastní analytický přínos.
  \item Procesní důkazy: drafty, commity, screencast, logy.
\end{enumerate}

\bigskip

\textbf{Doporučení k citování AI‑výstupů (stručně):}
\begin{itemize}
  \item APA: \textit{OpenAI. (2024). ChatGPT (verze X) [Large language model]. Výstup použit v úloze: popis a prompt.}
  \item IEEE: \textit{OpenAI, "ChatGPT (verze X)," 2024. Použito pro: [popis úlohy], prompt: "...".}
  \item Interní pole: \textit{Nástroj: ..., Prompt (zkráceně): "...", Vlastní přínos: "..."}
\end{itemize}

\bigskip

Poznámky z praxe:
\begin{itemize}
  \item U integrací v Microsoft 365 (např. Copilot) uveďte možnost generování výstupů do školních šablon \cite{kb_ai_101_tipu_pdf_04e4a115_page_44_2}.
  \item Ukázky AI ve výuce lze doplnit odkazy na Minecraft Education (Hour of Code) \cite{kb_ai_101_tipu_pdf_04e4a115_page_15_2}.
  \item U technických úloh požadujte procesní artefakty — inspirace: "Pokrok v matematice" \cite{kb_ai_101_tipu_pdf_04e4a115_page_8_1}.
\end{itemize}

Pro rychlé nasazení použijte tyto šablony jako výchozí bod a odkažte studenty na living document na ai.vut.cz; upravte text podle kategorie úloh (zakázáno / povoleno s deklarací / vyžadováno využití AI).